%! Special Directions --- Setting Up Eigenvalues — Unit 6, Lesson 6.0 — Exit Ticket (Answer Key)
\documentclass[10pt]{article}
\usepackage{linalg-article}
\usepackage{linalg-key}   % pulls in -boxes; do NOT also load -boxes
\newcommand{\TallMath}[1]{$\displaystyle #1\rule[-1.4em]{0pt}{3.2em}$}
\newcommand{\xx}{\mathbf{x}}

\begin{document}

\pageheader{Unit 6, Lesson 6.0}{Exit Ticket --- Answer Key}
\namedateperiod

\vspace{0.15in}

No notes. Show your reasoning. Use $A=\begin{bmatrix} 3 & 1 \\ 1 & 3 \end{bmatrix}$ for items 1--2.

\begin{enumerate}[leftmargin=1.8em, itemsep=16pt]
  \item \textbf{Test an eigenvector.} Compute $A\begin{bmatrix} 1 \\ 1 \end{bmatrix}$. Is
        $\begin{bmatrix} 1 \\ 1 \end{bmatrix}$ an eigenvector of $A$? If so, what is the eigenvalue
        $\lambda$?
        \ansline{$A\begin{bsmallmatrix} 1\\1 \end{bsmallmatrix}=\begin{bsmallmatrix} 4\\4 \end{bsmallmatrix}=4\begin{bsmallmatrix} 1\\1 \end{bsmallmatrix}$. Yes --- an eigenvector with $\lambda=4$.}
  \item \textbf{A vector that is \emph{not} special.} Compute $A\begin{bmatrix} 1 \\ 0 \end{bmatrix}$.
        Is $\begin{bmatrix} 1 \\ 0 \end{bmatrix}$ an eigenvector? How do you know?
        \ansline{$A\begin{bsmallmatrix} 1\\0 \end{bsmallmatrix}=\begin{bsmallmatrix} 3\\1 \end{bsmallmatrix}$, which is \emph{not} a multiple of $\begin{bsmallmatrix} 1\\0 \end{bsmallmatrix}$ --- so no, it is turned to a new direction.}
  \item \textbf{Explain.} In one sentence, what does it mean --- \emph{geometrically} --- for a nonzero
        vector $\xx$ to satisfy $A\xx=\lambda\xx$?
        \ansline{$A$ leaves $\xx$'s direction (its line) unchanged and only stretches it by the factor $\lambda$ (flipping it if $\lambda<0$) --- the matrix acts like the single number $\lambda$ along that direction.}
\end{enumerate}

\begin{teachernote}
Sort responses into ``got it / arithmetic slip / concept slip.'' Item~1 is a clean eigen-test
($\lambda=4$); item~2 checks that students actually verify ``multiple of $\xx$'' rather than assuming
every vector is special. Item~3 is the target sentence: an eigenvector is a direction the matrix only
\emph{stretches} (by $\lambda$). If many stumble on item~3, open Lesson~6.1 with a 60-second
``compute $A\xx$, is it a multiple of $\xx$?'' re-warm.
\end{teachernote}

\end{document}
